
\subsection{The whole septic geometric action}

Put $F=\mathbb Q(\zeta_7)$ and $F_0=\mathbb Q(\sqrt{-7})$,
and write $H=\operatorname{Gal}(F/F_0)=\{1,2,4\}$.
Let
\[
 A=\operatorname{Jac}(v^2+v=u^7),\qquad
 B=\operatorname{Jac}(y^7=x(1-x)^2),
\]
and let $E$ be the elliptic curve
$y^2+xy=x^3-x^2-2x-1$ of conductor $49$ and CM by
$\mathcal O_{F_0}$.  Its discriminant is $-7^3$ and
its $j$-invariant is $-3375$.
The substitution $(x_0,y_0)=(x-2,2y+x)$ gives
$y_0^2=4x_0^3+21x_0^2+28x_0$, the classical quotient in
\citet[Eq.~(2.10), pp.~68--69]{Elkies1998}.
All character comparisons below are made over $F$.
No rational splitting $B\sim E^3$ is asserted.

\begin{lemma}[Two distinct norm operations]
\label{lem:lm-septic-cm-comparison}
Let $\eta_A,\eta_B$ be the Jacobi Hecke characters over
$F$ for the triples $(1,1,5)$ and $(1,2,4)$, and let
$\eta_E$ be the CM Hecke character over $F_0$, choosing
the one of its two embeddings compatible with $H$.
Define
\[
 \eta_E^F(\mathfrak a)=
 \eta_E(N_{F/F_0}^{\mathrm{ideal}}\mathfrak a),\qquad
 \nu(\mathfrak a)=N_{F/\mathbb Q}\mathfrak a.
\]
Then, on ideals prime to $7$,
\[
 \eta_B=\eta_E^F,\qquad
 N_{F/F_0}^{\mathrm{coefficient}}(\eta_A)
       =\nu\,\eta_E^F.
 \tag{$\dagger$}\label{eq:lm-septic-character-identities}
\]
The superscripts distinguish the norm of the input ideal
from the norm of a character value.
\end{lemma}
\begin{proof}
The class number of $F$ is one.  Indeed, Minkowski's
bound is less than five, while primes above $2$ have
norm eight and primes above $3$ have norm $3^6$.
There is no nontrivial ideal of norm at most four.

For $A$, the reflex infinity type is
$R=\{1,4,5\}$; for $B$, it is $H=\{1,2,4\}$.
The latter is exactly the infinity type of $\eta_E^F$.
The CM character of $E/F_0$ has conductor the prime
above seven: the induction conductor formula is
$49=7\cdot7$.
Its local unit character is the nontrivial quadratic
residue character.  Pullback through the degree-three
totally ramified extension $F/F_0$ preserves that
quadratic unit character.

For $B$, \citet[Proposition~3.12(1), author p.~16]{ShuExpanded2021} applies because
\[
 1^1 2^2(4-7)^4=324,\qquad
 \frac{324^6-1}{7}\equiv0\pmod7.
\]
Thus $\eta_B$ has conductor exponent one at seven.
Equation~(3.14) expresses its unit character as a
seventh-power Hilbert-symbol character times the
quadratic character of conductor one.
The order-seven factor is trivial on units in the
conductor-one case.  Hence its unit character is
the same quadratic character as that of $\eta_E^F$.
Both characters are unramified away from seven.
Their quotient therefore has trivial infinity type
and is unramified at every finite place.  It is a
finite ideal-class character, so class number one
makes it trivial.  This proves the first identity.

For the second identity, the coefficient norm
multiplies the three coefficient conjugates indexed
by $H$.  Of the elements of $R$, two lie in $H$ and
one in its other coset.  Its infinity type is
therefore twice $H$ plus once $-H$.
Dividing by the ideal norm character $\nu$, whose
infinity type is once each of the six embeddings,
leaves exactly the type $H$ of $\eta_E^F$.

At seven, the finite unit character of $\eta_A$
has image $\mu_{14}$.  Coefficient norm to $F_0$
kills $\mu_7$, since $1+2+4=7$, and sends $-1$
to $-1$.  It therefore leaves precisely the same
quadratic unit character.  The character $\nu$ is
unramified on units.  Away from seven all characters
are unramified.  The quotient of the two sides of
the second identity is again an everywhere-unramified
finite ideal-class character of $F$, hence trivial.
\end{proof}

\begin{lemma}[The integral septic isogeny]\label{lem:septic-seven-isogeny}
The five cyclic quotient maps of $X_7$ induce homomorphisms
$\Phi:J_7\to A^3\times B^2$ and $\Psi:A^3\times B^2\to J_7$
such that $\Phi\Psi=[7]$ and $\Psi\Phi=[7]$. Consequently
\[
 T_3J_7\simeq(T_3A)^{\oplus3}\oplus(T_3B)^{\oplus2}
\]
integrally and $G_\Q$-equivariantly.
\end{lemma}
\begin{proof}
For $a=1,\ldots,5$, use the quotient curve
$C_a:y^7=x(1-x)^a$ and the rational morphism determined by
\[
 x=-X^7/Z^7,\qquad y=(-1)^{a+1}XY^a/Z^{a+1}.
\]
Its geometric deck subgroup in $\mu_7^2$ is
$H_a=\{(\zeta_7^u,\zeta_7^v):u+av=0\}$.
The thirty characters occurring in $H^1(X_7)$ are the pairs
$(i,j)$ with $i,j,i+j\ne0$ modulo seven. They partition into
the five lines $j=ai$ just listed. The triples for $a=1,3,5$
are equivalent to $(1,1,5)$ by coordinate permutations and
scaling all exponents by a unit modulo seven; those for
$a=2,4$ are equivalent to $(1,2,4)$. The resulting monomial
changes of cyclic coordinates are defined over $\Q$, so
the product of their Jacobians is $A^3\times B^2$.

Let $f_a$ denote these quotient maps. The trace identity gives
$f_a^*f_{a,*}=\sum_{h\in H_a}h$ on $J_7$ and
$f_{a,*}f_a^*=[7]$ on $\Jac(C_a)$. On a character line the
first operator is seven if the line is fixed by $H_a$ and
zero otherwise. Thus $f_{a,*}f_b^*=0$ for $a\ne b$, and
$\sum_a f_a^*f_{a,*}=[7]$. These identities on a rational
Tate module are identities of homomorphisms of abelian varieties,
by faithfulness of that realization. Taking the product of
the pushforwards and the sum of the pullbacks proves both
composition formulas. Since seven is a unit in $\Z_3$, they
give the asserted integral Tate-module isomorphism.
\end{proof}

\begin{proposition}[All compact residual factors are controlled]
\label{prop:lm-whole-septic-action}
Let $K=\mathbb Q(A[3])$.  Then $E[3]$ is rational over
$K$, and the image of $G_K$ on the entire rational
Tate module of the degree-seven Fermat Jacobian is
pro-$3$.  Consequently, with a standard unit tangent,
the actual action of $G_K$ on the entire geometric
pro-$3$ group $N_7$ has pro-$3$ image.
The same is true for $N_{7\cdot3^r}$, $r\ge0$.
All these actions are unramified outside $\{3,7\}$.
\end{proposition}
\begin{proof}
On $G_F$, the second identity in
\eqref{eq:lm-septic-character-identities} gives the
$3$-adic character identity
\[
 \rho_{E,3}=
 \chi_3^{-1}N_{F_3/(F_0)_3}^{\mathrm{coefficient}}(\rho_{A,3}),
\]
up to the harmless choice of the conjugate CM embedding.
The field $K$ contains $F$ and $\mu_3$.
On $G_K$, both $\rho_{A,3}$ and $\chi_3$ are congruent
to one modulo three, so the same is true of
$\rho_{E,3}$.  Thus $E[3]\subset E(K)$.
The images on $T_3A$ and $T_3E$ over $G_K$ are pro-$3$.

Lemma~\ref{lem:septic-seven-isogeny} gives the decomposition
of $V_3J_7$ into three copies of $V_3A$ and two of $V_3B$.
After restriction to $G_F$, the first character
identity in \eqref{eq:lm-septic-character-identities}
identifies $V_3B$ with three copies of $V_3E$:
the six coefficient embeddings yield three copies
of each of the two $F_0$-embeddings.
Hence $V_3J_7|_{G_K}$ has pro-$3$ image.
This topological property does not depend on the
choice of a stable integral lattice.

The boundary of the Fermat curve consists of three
copies of $\mu_7$.  Since $K$ contains both $\mu_7$
and $\mu_3$, the boundary modulo three is trivial
over $G_K$.  The compact and boundary terms show
that the image on $N_7/\Phi(N_7)$ is a three-group.
The automorphism kernel over the Frattini quotient
is pro-$3$, proving the assertion for the full $N_7$.

Finally $U_{7\cdot3^r}\to U_7$ is a genuine geometric
three-power cover.  Compatible unit tangents,
$u_7=u_{7\cdot3^r}^{3^r}$, identify the upstairs group
with an arithmetic-stable open subgroup, so the
action statement passes upstairs.  Good logarithmic
reduction of the pair and the extending tangent
give unramified actual action outside $\{3,7\}$.
\end{proof}

\begin{corollary}[A finite sufficient gate for $7\cdot3^r$]
\label{cor:lm-septic-class-gate}
If $\operatorname{Cl}(\mathbb Q(A_7[3]))[3]=0$, then
for every $r\ge1$ all full, abelian, finite lower-central,
and finite-derived geometric pro-$3$ Selmer sections
of $U_{7\cdot3^r}/\mathbb Q$ are globally cuspidal,
with the three rational cusp packets as before.
The same conclusion holds for an arithmetic-stable
quotient retaining abelianization and having finite
free ordinary lower-central factors. Moreover the compact
integral Selmer localization at three is injective for each such $n$.
\end{corollary}
\begin{proof}
The ray-class lemma makes the local map at three
surjective onto the maximal $\{3,7\}$-ramified
pro-$3$ arithmetic source over $K$; the residual
decomposition group already equals
$\operatorname{Gal}(K/\mathbb Q)$.
The preceding proposition places every compact,
boundary, and central coefficient action in that
source.  The cocycle-graph argument therefore gives
the required integral cohomology localization
injections.  The cubic compact detector, exact
three-primary order of cusp differences, and the
same central induction as for $5\cdot3^r$ apply.
More explicitly, the coefficient-graph homomorphism proves
$H^1(G_{\Q,\{3,7\}},L)\hookrightarrow H^1(\Q_3,L)$
for each finite free coefficient module just described, exactly as
in Lemma~\ref{lem:lm-five-power-action}. Compact Selmer classes
are unramified outside $\{3,7\}$ by good reduction, proving
the compact assertion. Positive good-prime weights give zero
global invariants in every lower-central factor; source surjectivity
gives the same at three. Thus all hypotheses of
Proposition~\ref{prop:lm-real-mixed-higher-criterion} hold with $q=7$.
The class-group hypothesis is verified below in
Proposition~\ref{prop:septic-three-class}, using the explicit field
calculation and degree-$56$ reduction. This corollary then supplies
the septic case of Theorem~\ref{thm:unrestrictedthreefamiliesintro}.
\end{proof}
